\documentclass[10pt,a4paper]{amsart}
\usepackage[margin=19mm]{geometry}
\usepackage{amsmath,amssymb,mathtools}
\usepackage[scr=rsfs,cal=euler]{mathalfa}
\usepackage{xfrac}
\usepackage[unicode,hidelinks,bookmarks=false]{hyperref}
\newcommand{\ie}{i.e.\ }
\newcommand{\cf}{cf.\ }
\newcommand{\resp}{respectively\ }
\newcommand{\Subsection}{\subsection}
\providecommand{\nobreakdash}{\nobreak\hbox{-}}
\setlength{\emergencystretch}{2em}
\multlinegap0pt
\usepackage{amssymb}
\usepackage{mathtools}
\newtheorem{theorem}{Theorem}
\newtheorem{lemma}{Lemma}
\newtheorem{proposition}{Proposition}
\newcommand{\Q}{\mathbb{Q}}
\newcommand{\R}{\mathbb{R}}
\newcommand{\Z}{\mathbb{Z}}
\DeclareMathOperator{\cone}{cone}
\newcommand{\lam}{\lambda}
\title{Positivity of stretched Littlewood--Richardson\\ coefficients for partitions of length at most five}
\author{Alper Ferudun}
\date{Source: arXiv:2607.22301v2 (CC BY 4.0)}
\begin{document}
\maketitle

\noindent\emph{Source and licence.} Positivity of stretched Littlewood--Richardson\\ coefficients for partitions of length at most five. Version arXiv:2607.22301v2 (CC BY 4.0), distributed by arXiv under the Creative Commons Attribution 4.0 International licence, http://creativecommons.org/licenses/by/4.0/, as recorded in the arXiv licence metadata for this exact version and shown on the abstract page https://arxiv.org/abs/2607.22301v2. Full licence notice: \texttt{licenses/arxiv-2607.22301-CC-BY-4.0.txt}.\par\smallskip

\noindent\textit{Editorial scope.} Counted unit: the article's proposition prop:certificate (source0.tex lines 565--575). The article's complete original proof of it is delivered with it (source0.tex lines 592--606, one \texttt{proof} environment of 15 lines, which the article places away from the statement). No mathematics is written, repaired or re-derived: the statement and the proof are the article's own lines, in the article's own notation, and no retained line is substituted or re-flowed.  Retained uncounted as the source-local context the proof needs: lines 136--139; lines 260--265; lines 536--539.  Nothing was omitted from any retained span, so the package carries no omission record.  No retained line contains a citation, so the wrapper carries no bibliography and no bibliography entry.  No retained line contains a figure environment, an image-inclusion command or drawing code, so no image is delivered and the package declares no asset dependency.  Editorial formatting only, in addition: an AMS \texttt{amsart} preamble that restates the article's own declarations and macros used by the retained lines, namely the environments \texttt{theorem}, \texttt{lemma}, \texttt{proposition} on separate counters, together with the standard packages those lines need.  The retained environments print in this document as Theorem 1, Lemma 1, Proposition 1 in order of appearance, whereas the article prints the same items with the numbering of its own full document.  The complete native source of the article as distributed in the arXiv e-print for this exact version is bundled unchanged as \texttt{sources/205929/source0.tex}.
\begin{theorem}\label{thm:linear}
For every full-dimensional five-part hive polytope the linear Ehrhart
coefficient is nonnegative.
\end{theorem}

\begin{lemma}\label{lem:dilation}
Fix $r$ and let $N_r$ be the (finite) set of primitive rhombus normals. Suppose
that for every \emph{lattice} polytope $P\subset\R^{D}$ whose facet normals lie
in $N_r$ one has $[t^{k}]L_{P}(t)\ge 0$. Then $[t^{k}]P^{\nu}_{\lam\mu}(t)\ge 0$
for all partitions of length at most $r$.
\end{lemma}

\begin{equation}\label{eq:corrected}
  a_1(H)=\sum_{\text{edges }E}\ell_{\Z}(E)\Bigl[\alpha\bigl(N(H,E)\bigr)
  +\sum_{F\supset E,\ \dim F=2}\bigl\langle y_{N(H,F)},u_{E/F}\bigr\rangle\Bigr],
\end{equation}

\begin{proposition}\label{prop:certificate}
Let $\Sigma_{\mathcal T}$ be the set of all closed pointed rank-five types
$\sigma$ (over all edge directions) having at least one facet in $\mathcal T$;
$|\Sigma_{\mathcal T}|=87{,}127$, over $551$ directions. There are rational
vectors $y_\tau\in\Q^{2}$, $\tau\in\mathcal T$, with denominators at most
$10^{4}$, such that for every $\sigma\in\Sigma_{\mathcal T}$
\[
  \alpha(\cone\sigma)+\sum_{\tau\in\mathcal T\text{ facet of }\sigma}
  \langle y_\tau,u_{\sigma/\tau}\rangle\ \ge\ 1.20\times10^{-3}.
\]
\end{proposition}

\begin{proof}[Proof of Theorem~\ref{thm:linear}]
Let $H$ be a full-dimensional five-part hive and, for each two-face $F$, let
$y_{N(H,F)}$ be the vector of Proposition~\ref{prop:certificate} if the type of
$F$ lies in $\mathcal T$ and $0$ otherwise. By \eqref{eq:corrected} it suffices
to show that every bracket is nonnegative. Let $E$ be an edge of type $\sigma$.
If some facet of $\sigma$ lies in $\mathcal T$ then $\sigma\in\Sigma_{\mathcal T}$
and the bracket is at least $1.20\times10^{-3}$ by
Proposition~\ref{prop:certificate}. Otherwise the bracket equals
$\alpha(\cone\sigma)$. If the pulling triangulation of $\sigma$ contains a
negative cell then $\sigma\in\Sigma^{-}$; it cannot lie in $\Sigma^{-}_{<0}$,
because every facet of such a cone is in $\mathcal T$, so
$\alpha(\cone\sigma)\ge0$. If it contains no negative cell then
$\alpha(\cone\sigma)$ is a sum of positive cell weights. Hence $a_1(H)\ge0$.
Rational hives are covered by Lemma~\ref{lem:dilation}.
\end{proof}

\end{document}
